\begin{tikzpicture}[
  font=\scriptsize,
  dalam/.style={draw, rounded corners, fill=praditagreen!15, text width=1.6cm,
                minimum height=0.8cm, align=center},
  luar/.style={draw, rounded corners, fill=blue!8, text width=2.05cm,
               minimum height=0.8cm, align=center},
  manual/.style={draw, rounded corners, fill=red!22, text width=1.05cm,
                 minimum height=0.5cm, align=center, font=\tiny, inner sep=1.5pt},
  gerbang/.style={draw, rounded corners, fill=orange!12, text width=1.05cm,
                  minimum height=0.5cm, align=center, font=\tiny\bfseries, inner sep=1.5pt},
  batas/.style={draw=praditagreen, dashed, thick, rounded corners, inner sep=4pt},
  judul/.style={font=\scriptsize\bfseries, align=center},
  ket/.style={font=\tiny, align=center},
  panah/.style={-{Stealth[length=2mm]}, thick}
]
  % Panel kiri: tiap mitra melintasi batas dengan caranya sendiri.
  \node[judul] at (2.8,1.75) {Sebelum: tiap mitra punya cara sendiri};
  \node[dalam] (a1) at (0.9,1.0) {Layanan klinis};
  \node[dalam] (a2) at (0.9,0) {Data pasien};
  \node[dalam] (a3) at (0.9,-1.0) {Tagihan dan rekap};
  \node[batas, fit=(a1)(a2)(a3)] (ba) {};
  \node[luar] (l1) at (4.8,1.0) {3 laboratorium\\mitra};
  \node[luar] (l2) at (4.8,0) {9 asuransi\\swasta};
  \node[luar] (l3) at (4.8,-1.0) {36 perusahaan\\klien};
  \draw[panah] (l1.west) -- node[manual] {surel} (l1.west -| ba.east);
  \draw[panah] (ba.east |- l2.west) -- node[manual] {9 portal} (l2.west);
  \draw[panah] (ba.east |- l3.west) -- node[manual] {\textit{spreadsheet}} (l3.west);
  \node[ket] at (2.8,-1.95) {Misalkan 259 jam kerja manual per bulan};

  % Panel kanan: batas tetap ada, tetapi cara melintasinya seragam dan terkelola.
  \begin{scope}[xshift=6.35cm]
    \node[judul] at (2.8,1.75) {Sesudah: batas punya gerbang terkelola};
    \node[dalam] (b1) at (0.9,1.0) {Layanan klinis};
    \node[dalam] (b2) at (0.9,0) {Data pasien};
    \node[dalam] (b3) at (0.9,-1.0) {Tagihan dan rekap};
    \node[batas, fit=(b1)(b2)(b3)] (bb) {};
    \node[luar] (r1) at (4.8,1.0) {3 laboratorium\\mitra};
    \node[luar] (r2) at (4.8,0) {9 asuransi\\swasta};
    \node[luar] (r3) at (4.8,-1.0) {36 perusahaan\\klien};
    \draw[panah] (r1.west) -- node[gerbang] {API hasil} (r1.west -| bb.east);
    \draw[panah] (bb.east |- r2.west) -- node[gerbang] {API klaim} (r2.west);
    \draw[panah] (bb.east |- r3.west) -- node[gerbang] {API rekap} (r3.west);
    \node[ket] at (2.8,-1.95) {Hak akses tetap diatur, biaya melintasi batas turun};
  \end{scope}
\end{tikzpicture}
